\section{结果}

\subsection{数据预处理}

在实验数据的预处理阶段，为确保分析结果的可靠性和有效性，本研究依据反应时异常（超过被试内均值$\pm 3$倍标准差）等标准，对原始数据进行了筛选。分析对象为正式 block 中的摆球刺激试次（\texttt{pendulumStimulus} + \texttt{block}），不含练习段。数据合并、质控与统计均使用 Python（pandas、pingouin、statsmodels）完成。

原始数据共 1248 个试次（8 名被试）。正式 block 试次 1000 个；经反应时质控剔除 20 个试次后，最终保留 980 个试次进入描述性与推断统计分析。各被试正式试次数均为 125（剔除前），25 个能量水平各 40 试次，无缺失值。被试 0010 的绝对误差均值（约 $34.77^\circ$）高于其余被试，已在质控报告中标记；主分析保留该被试，并在推断统计中报告剔除后的敏感性分析。

\subsection{描述性统计}

表~\ref{tab:desc-ebin} 记录了低、中、高三档能量条件下的有符号角偏差与绝对角误差。总体来看，绝对角误差随能量升高而增大：低能量档 $M = 9.38^\circ$（$SD = 10.03^\circ$），中能量档 $M = 16.75^\circ$（$SD = 16.11^\circ$），高能量档 $M = 20.62^\circ$（$SD = 19.47^\circ$）。中档较低档平均高出约 $7.36^\circ$，高档较低档平均高出约 $11.24^\circ$。

有符号偏差方面，三档能量条件下均值均接近 0（低：$M = 0.72^\circ$；中：$M = 1.22^\circ$；高：$M = -1.62^\circ$），提示整体上未表现出单一方向的一致性系统偏差。表~\ref{tab:desc-overall} 汇总了正式试次各变量的总体描述统计；表~\ref{tab:desc-subject} 显示被试间差异，其中被试 0010 的 $|{\Delta\theta}|$ 明显高于其他被试。

\begin{table}[htbp]
  \centering
  \caption{不同能量档下的角估计误差（M$\pm$SD）}
  \label{tab:desc-ebin}
  \begin{tabular}{lccc}
    \toprule
    能量档 & 有符号偏差 $\Delta\theta$ (°) & 绝对误差 $|\Delta\theta|$ (°) & $n$ \\
    \midrule
    低 & 0.72$\pm$13.72 & 9.38$\pm$10.03 & 355 \\
    中 & 1.22$\pm$23.22 & 16.75$\pm$16.11 & 352 \\
    高 & -1.62$\pm$28.34 & 20.62$\pm$19.47 & 273 \\
    \bottomrule
  \end{tabular}
\end{table}

\begin{table}[htbp]
  \centering
  \caption{正式 block 试次总体描述统计（M$\pm$SD）}
  \label{tab:desc-overall}
  \begin{tabular}{lc}
    \toprule
    变量 & M$\pm$SD \\
    \midrule
    机械能 $E$ (J) & 38.61$\pm$21.22 \\
    有符号偏差 $\Delta\theta$ (°) & 0.25$\pm$22.04 \\
    绝对误差 $|\Delta\theta|$ (°) & 15.16$\pm$16.00 \\
    反应时 RT (s) & 2.99$\pm$1.91 \\
    \bottomrule
  \end{tabular}
\end{table}

\begin{table}[htbp]
  \centering
  \caption{分被试绝对角误差与反应时（M$\pm$SD）}
  \label{tab:desc-subject}
  \begin{tabular}{lccc}
    \toprule
    被试 & $n$ & $|\Delta\theta|$ (°) & RT (s) \\
    \midrule
    0002 & 122 & 12.93$\pm$10.47 & 2.03$\pm$0.51 \\
    0003 & 124 & 14.68$\pm$12.62 & 2.22$\pm$0.89 \\
    0004 & 123 & 7.44$\pm$7.52 & 2.44$\pm$0.65 \\
    0006 & 122 & 14.39$\pm$11.53 & 3.81$\pm$1.54 \\
    0007 & 123 & 15.23$\pm$15.08 & 1.89$\pm$0.75 \\
    0008 & 122 & 13.07$\pm$12.50 & 1.79$\pm$0.66 \\
    0009 & 123 & 9.29$\pm$8.33 & 6.24$\pm$2.79 \\
    0010 & 121 & 34.50$\pm$26.08 & 3.51$\pm$0.69 \\
    \bottomrule
  \end{tabular}
\end{table}


图~\ref{fig:energy-residual} 展示了机械能 $E$ 与有符号角残差 $\Delta\theta$ 的散点分布及分箱中位趋势线，可见高能量试次多分布于较大残差区域。图~\ref{fig:energy-abs} 显示 $E$ 与 $|{\Delta\theta}|$ 呈正相关趋势。图~\ref{fig:residual-dist} 表明 $|{\Delta\theta}|$ 分布右偏，中位数（$Md = 10.24^\circ$）低于均值（$M = 15.16^\circ$）。图~\ref{fig:ebin-violin} 与图~\ref{fig:subject-box} 分别呈现三档能量与被试间的绝对误差分布；图~\ref{fig:subspace} 在 $(\theta,\,|\omega|/\omega_0)$ 相空间上以颜色编码 $|{\Delta\theta}|$，显示误差在相空间各区域均有分布。

\begin{figure}[htbp]
  \centering
  \includegraphics[width=0.85\linewidth]{energy_vs_residual_formal.png}
  \caption{正式 block：机械能 vs 有符号角残差}
  \label{fig:energy-residual}
\end{figure}

\begin{figure}[htbp]
  \centering
  \includegraphics[width=0.85\linewidth]{energy_vs_abs_error_formal.png}
  \caption{正式 block：机械能 vs 绝对角误差}
  \label{fig:energy-abs}
\end{figure}

\begin{figure}[htbp]
  \centering
  \includegraphics[width=0.9\linewidth]{residual_distribution_formal.png}
  \caption{正式 block：估计误差分布}
  \label{fig:residual-dist}
\end{figure}

\begin{figure}[htbp]
  \centering
  \includegraphics[width=0.75\linewidth]{e_bin_violin.png}
  \caption{三档能量下的绝对角误差分布}
  \label{fig:ebin-violin}
\end{figure}

\begin{figure}[htbp]
  \centering
  \includegraphics[width=0.85\linewidth]{subject_boxplot.png}
  \caption{被试间绝对角误差比较}
  \label{fig:subject-box}
\end{figure}

\begin{figure}[htbp]
  \centering
  \includegraphics[width=0.85\linewidth]{viz_subspace_error_human.png}
  \caption{相空间上的绝对角误差（人类数据）}
  \label{fig:subspace}
\end{figure}

各能量水平的细粒度描述统计见表~\ref{tab:desc-elevel}（\texttt{longtable}）。

\begin{longtable}{lcccc}
  \caption{各能量水平下的角估计误差（M$\pm$SD）} \label{tab:desc-elevel} \\
  \toprule
  $E$ (J) & $n$ & $\Delta\theta$ (°) & $|\Delta\theta|$ (°) \\
  \midrule
  \endfirsthead
  \multicolumn{4}{c}{\tablename\ \thetable{} -- 续表} \\
  \toprule
  $E$ (J) & $n$ & $\Delta\theta$ (°) & $|\Delta\theta|$ (°) \\
  \midrule
  \endhead
  \bottomrule
  \endfoot
  3.43 & 39 & 1.99$\pm$6.54 & 4.42$\pm$5.17 \\
  6.37 & 40 & -0.38$\pm$8.05 & 5.97$\pm$5.33 \\
  9.31 & 40 & 0.64$\pm$11.29 & 7.52$\pm$8.36 \\
  12.25 & 39 & 2.21$\pm$14.57 & 9.52$\pm$11.16 \\
  15.19 & 39 & 2.68$\pm$13.34 & 9.65$\pm$9.47 \\
  18.13 & 40 & 2.50$\pm$10.51 & 8.79$\pm$6.14 \\
  21.07 & 38 & -0.78$\pm$21.59 & 16.08$\pm$14.18 \\
  24.01 & 40 & -2.35$\pm$12.06 & 9.73$\pm$7.34 \\
  26.95 & 40 & 0.05$\pm$19.31 & 12.99$\pm$14.14 \\
  29.89 & 39 & 2.24$\pm$26.61 & 19.33$\pm$18.15 \\
  32.83 & 40 & -7.19$\pm$20.72 & 15.89$\pm$14.94 \\
  35.77 & 39 & 1.33$\pm$17.12 & 12.53$\pm$11.56 \\
  38.71 & 40 & 2.16$\pm$20.94 & 16.26$\pm$13.12 \\
  41.65 & 38 & 6.93$\pm$21.64 & 16.41$\pm$15.53 \\
  44.59 & 39 & 1.47$\pm$17.34 & 13.62$\pm$10.60 \\
  47.53 & 39 & 7.43$\pm$20.85 & 16.77$\pm$14.24 \\
  50.47 & 39 & 4.42$\pm$31.89 & 21.70$\pm$23.53 \\
  53.41 & 39 & -7.46$\pm$25.35 & 18.21$\pm$18.96 \\
  56.35 & 39 & 4.14$\pm$24.50 & 17.31$\pm$17.62 \\
  59.29 & 40 & -3.62$\pm$19.69 & 14.47$\pm$13.66 \\
  62.23 & 38 & 1.91$\pm$27.95 & 20.87$\pm$18.37 \\
  65.17 & 38 & -12.40$\pm$33.67 & 25.00$\pm$25.50 \\
  68.11 & 40 & -2.63$\pm$25.45 & 20.54$\pm$14.91 \\
  71.05 & 38 & -7.72$\pm$25.40 & 21.39$\pm$15.37 \\
  73.99 & 40 & 8.47$\pm$35.34 & 24.94$\pm$26.16 \\
\end{longtable}


\subsection{推断统计}

对绝对角误差进行单因素（能量档：低/中/高）重复测量方差分析，以被试在各档内的试次均值为分析单元（$n = 8$）。Mauchly 球形性检验显著（$W = 0.226$，$p = 0.012$），故参考 Greenhouse--Geisser 校正结果：能量档主效应显著（$F(2, 14) = 41.132$，$p < .001$，$\eta_p^2 = 0.855$），表明机械能水平对摆角估计精度具有显著影响。未校正 $p$ 值同样显著（$p < .001$).

事后检验（Bonferroni 校正）显示：低与中比较（Bonferroni 校正）达到显著（$t(7) = -4.756$，$p_{\mathrm{corr}} = 0.006$，Cohen's $d = -0.916$）。 低与高比较（Bonferroni 校正）达到显著（$t(7) = -7.722$，$p_{\mathrm{corr}} = 0.000$，Cohen's $d = -1.427$）。 中与高比较（Bonferroni 校正）达到显著（$t(7) = -8.420$，$p_{\mathrm{corr}} = 0.000$，Cohen's $d = -0.391$）。

\begin{table}[htbp]
  \centering
  \caption{能量档（低/中/高）对绝对角误差的重复测量方差分析}
  \label{tab:infer-rm-anova}
  \begin{tabular}{lccc}
    \toprule
    因素 & $F$ & $p$ & $\eta_p^2$ \\
    \midrule
    能量档 & 41.132 & < .001 & 0.855 \\
    \bottomrule
  \end{tabular}
  \begin{flushleft}
    \footnotesize
    注：Mauchly 球形性检验 $W$ = 0.226，$p$ = = .012。
  \end{flushleft}
\end{table}

\begin{table}[htbp]
  \centering
  \caption{能量档两两比较的事后检验（Bonferroni 校正）}
  \label{tab:infer-posthoc}
  \begin{tabular}{lccc}
    \toprule
    比较 & $t$ & $p_{\mathrm{corr}}$ & Cohen's $d$ \\
    \midrule
    低 vs 中 & -4.756 & = .006 & -0.916 \\
    低 vs 高 & -7.722 & < .001 & -1.427 \\
    中 vs 高 & -8.420 & < .001 & -0.391 \\
    \bottomrule
  \end{tabular}
\end{table}


为进一步检验系统性有符号偏差，对 $\Delta\theta$ 进行单样本 $t$ 检验（检验值 $= 0^\circ$）。全体试次均值（$M = 0.25^\circ$，$SD = 22.04^\circ$）未显著偏离 0（$t(979) = 0.356$，$p = 0.722$，Cohen's $d = 0.011$），各能量档亦均未表现出显著单向偏差（见表~\ref{tab:infer-ttest}）。

\begin{table}[htbp]
  \centering
  \caption{有符号角偏差 $\Delta\theta$ 的单样本 $t$ 检验（检验值 = 0°）}
  \label{tab:infer-ttest}
  \begin{tabular}{lcccccc}
    \toprule
    条件 & $M$ (°) & $SD$ (°) & $t$ & $df$ & $p$ & Cohen's $d$ \\
    \midrule
    全体 & 0.25 & 22.04 & 0.356 & 979 & = .722 & 0.011 \\
    低 & 0.72 & 13.72 & 0.994 & 354 & = .321 & 0.053 \\
    中 & 1.22 & 23.22 & 0.988 & 351 & = .324 & 0.053 \\
    高 & -1.62 & 28.34 & -0.944 & 272 & = .346 & -0.057 \\
    \bottomrule
  \end{tabular}
  \begin{flushleft}
    \footnotesize 注：$p$ 值采用双尾检验。
  \end{flushleft}
\end{table}


Pearson 相关分析表明，试次级 $E$ 与 $|{\Delta\theta}|$ 显著正相关（$r = 0.309$，$p < .001$）；被试内中心化后相关仍显著（$r = 0.356$，$p < .001$），支持能量升高伴随估计难度增加的探索性结论。

作为敏感性分析，线性混合模型（$|{\Delta\theta}| \sim E_{z} + \mathrm{show\_T}_{z} + (1|\mathrm{subject})$）显示：标准化能量 $E_z$ 的固定效应显著（$\beta = 4.980$，$p < .001$），可见时长 $\mathrm{show\_T}_z$ 效应不显著（$p = 0.883$）。剔除被试 0010 后（$n = 859$ 试次），$E_z$ 效应仍显著（$\beta = 4.323$，$p < .001$），结果方向保持一致。

鉴于本研究被试量较小（$n = 8$），上述推断统计应视为\textbf{探索性分析}，结论侧重效应方向与描述性模式，有待更大样本进一步验证。
